\section{Some True/False Questions}

\textbf{(a) False.}
Let
\[
    A =
    \begin{bmatrix}1&0\\0&1\end{bmatrix},
    \qquad
    B =
    \begin{bmatrix}0&1\\1&0\end{bmatrix}.
\]
Both have nonnegative entries and both are onto, but
\[
    A+B =
    \begin{bmatrix}1&1\\1&1\end{bmatrix}
\]
has rank one, so it is not onto.

\textbf{(b) True.}
If
\[
    Ax=0,\qquad (A+B)x=0,\qquad (A+B+C)x=0,
\]
then \(Bx=0\) follows from the first two equations, and \(Cx=0\)
follows from all three. Thus \(x\in\operatorname{null}(A)\cap
\operatorname{null}(B)\cap\operatorname{null}(C)\).% The reverse
%containment is immediate.

\textbf{(c) False.}
As a counterexample, taake \(A=0\), \(B=1\), and \(C=1\). Then
\[
    \operatorname{null}
    \left(
    \begin{bmatrix}
        A\\ AB\\ ABC
    \end{bmatrix}
    \right)
    =
    \mathbf{R},
\]
but
\[
    \operatorname{null}(A)\cap\operatorname{null}(B)\cap
    \operatorname{null}(C)
    =
    \{0\}.
\]

\textbf{(d) True.}
For any \(x\),
\[
    x^T(B^T A^T A B+B^T B)x
    =
    \|ABx\|^2+\|Bx\|^2.
\]
This is zero if and only if \(Bx=0\). Hence the nullspace is
\(\operatorname{null}(B)\).

\textbf{(e) True.}
The rank of a block diagonal matrix is the sum of the ranks of its
diagonal blocks:
\[
    \operatorname{rank}
    \begin{bmatrix}
        A&0\\0&B
    \end{bmatrix}
    =
    \operatorname{rank}(A)+\operatorname{rank}(B).
\]
If the block diagonal matrix is full rank, then neither diagonal block
can be rank deficient. Therefore both \(A\) and \(B\) are full rank.

\textbf{(f) True.}
The matrix \(\begin{bmatrix} A&0\end{bmatrix}\) has the same column space
as \(A\). If \(\begin{bmatrix} A&0\end{bmatrix}\) is onto, then \(A\) is
onto, so \(A\) is full rank.

\textbf{(g) True.}
Since \(A^2\) is defined, \(A\) is square. If \(A^2\) is onto, then for
every \(y\) there is some \(x\) such that
\[
    A^2x = y.
\]
Let \(z=Ax\). Then
\[
    Az = y.
\]
So every \(y\) is in the range of \(A\), which means \(A\) is onto.

\textbf{(h) False.}
Let
\[
    A =
    \begin{bmatrix}
        1\\0
    \end{bmatrix}.
\]
Then \(A^T A = [1]\), which is onto as a map \(\mathbf{R}\to\mathbf{R}\).
But \(A:\mathbf{R}\to\mathbf{R}^2\) is not onto.

\textbf{(i) True.}
Suppose \(\alpha_i\geq 0\) and
\[
    \sum_{i=1}^k \alpha_i u_i=0.
\]
Then
\[
    0 =
    \left\|\sum_{i=1}^k \alpha_i u_i\right\|^2
    =
    \sum_{i=1}^k\sum_{j=1}^k
    \alpha_i\alpha_j u_i^T u_j.
\]
Every term is nonnegative, and the diagonal terms are
\(\alpha_i^2\|u_i\|^2\). Since each \(u_i\neq 0\), all \(\alpha_i\) must
be zero.

\textbf{(j) True.}
If
\[
    Ax+By=0,
\]
then \(Ax=-By\). The vector \(Ax\) lies in the column space of \(A\), and
\(By\) lies in the column space of \(B\). Since \(A^T B=0\), these two
column spaces are orthogonal. Hence
\[
    0=\|Ax+By\|^2=\|Ax\|^2+\|By\|^2,
\]
so \(Ax=0\) and \(By=0\). Because \(A\) and \(B\) are skinny and full
rank, \(x=0\) and \(y=0\). Therefore \([A\ B]\) has independent columns,
so it is skinny and full rank.
